\chapter{Обзор области исследований}\label{ch:ch1}
Глава 1 посвящена обзору области исследований, связанной с оптимизацией SQL-запросов. В ней рассматриваются историческое развитие методов планирования, устройство современных оптимизаторов, ограничения классических стоимостных подходов и переход к методам оптимизации запросов с использованием машинного обучения и генеративных моделей \cite{10.1145/275487.275492, Lan_2021, ZOU2022100304, 10.1145/3626246.3654692}. Отдельно обсуждаются вопросы экспериментального тестирования и выбора бенчмарков, применяемых для оценки методов оптимизации запросов.

\section{Историческое развитие алгоритмов планирования запросов}

Проблема автоматического выбора плана выполнения возникла вместе с реляционной моделью данных. В навигационных СУБД разработчик фактически сам задавал путь доступа к данным и порядок обработки записей. Переход к декларативной реляционной модели привёл к необходимости автоматически определять эффективную стратегию выполнения запроса \cite{10.1145/362384.362685, 10.1145/356662.356666, 10.1145/800227.806919, 10.1145/960124.806919}.

Ключевым этапом стало появление оптимизатора System R, описанного в классических работах IBM \cite{10.1145/582095.582099, 10.1145/320455.320457, 10.1145/358769.358784}. Именно там были сформулированы базовые принципы, определившие развитие области на десятилетия вперёд: использование эквивалентных преобразований запроса, применение стоимостной модели для сравнения альтернатив и перебор пространства планов с помощью динамического программирования \cite{10.1145/582095.582099, 10.1145/320455.320457, 10.1145/275487.275492}.

В дальнейшем развитие шло по двум основным направлениям. С одной стороны, совершенствовались архитектуры оптимизаторов. В работах, связанных с Volcano и Cascades, были предложены расширяемые фреймворки, в которых правила преобразования и физические свойства плана оформляются единообразно, а эквивалентные выражения хранятся в общей структуре memo \cite{10.1109/69.273032, 344061, DBLP:journals/debu/Graefe95a}. Это упростило поддержку новых операторов и сделало оптимизаторы более модульными \cite{10.1145/152610.152611, 10.1145/275487.275492}.

С другой стороны, стало очевидно, что даже при удачной архитектуре пространство планов растёт слишком быстро. По мере увеличения числа соединений полный перебор перестаёт быть практически выполнимым, поэтому в оптимизаторах начали широко использоваться эвристики, отсечения и приближённые методы поиска, включая стохастические алгоритмы \cite{10.1145/93605.98740, 10.1007/s007780050040, 10.1007/s00778-017-0480-7}.

К рубежу 2000-х годов сложилась устойчивая парадигма классического оптимизатора: правила переписывания запроса, стоимостная модель и ограниченный перебор альтернатив \cite{10.1145/275487.275492, 10.1145/152610.152611}. При всех успехах этой схемы стало ясно, что её качество сильно зависит от точности статистических оценок, адекватности стоимостной модели и способности учитывать реальные условия выполнения \cite{10.1007/s00778-017-0480-7, Lan_2021}.

Следующим этапом стало появление методов адаптации и обучения. Ранние работы, такие как IBM LEO, использовали обратную связь от выполненных запросов для уточнения внутренних оценок \cite{10.1147/sj.421.0098, 10.5555/645927.672349}. Параллельно появились исследования, в которых машинное обучение применялось для предсказания времени выполнения или коррекции отдельных компонентов оптимизатора. Например, в работе Akdere et al.\cite{10.1109/ICDE.2012.64} показано, что характеристики плана могут использоваться для достаточно точного прогноза фактического времени выполнения запроса.

К 2010-м годам рост объёмов данных, усложнение рабочих нагрузок и прогресс в обучении сложных моделей привели к формированию нового поколения исследований, в которых обучаемые методы стали рассматриваться как средство улучшения отдельных компонентов оптимизатора \cite{Lan_2021, ZOU2022100304, 10.1145/3626246.3654692}.

\section{Оптимизатор SQL-запросов в СУБД}

Современные реляционные СУБД используют оптимизатор запросов, или оптимизатор, для преобразования декларативного SQL-запроса в план выполнения~--- последовательность физических операций над данными. Оптимизатор является частью процессора запросов СУБД и получает на входе внутреннее представление запроса после синтаксического разбора и логического переписывания. Его задача состоит в том, чтобы выбрать наиболее эффективный способ выполнения запроса с точки зрения ожидаемой производительности. 

Рисунок \ref{fig:postgresql_planner} показывает место оптимизатора в общей архитектуре СУБД на примере PostgreSQL. После разбора запроса и применения правил переписывания оптимизатор формирует дерево плана выполнения, включающее конкретные физические операции: способы сканирования таблиц, алгоритмы соединения, сортировку, фильтрацию, агрегирование и другие шаги. При выборе плана оптимизатор использует не только переписанное представление запроса, но и информацию о схеме базы данных, включая определения отношений, индексов и ограничений, а также статистическую информацию о данных. Сформированный план затем передаётся исполнителю запросов.

\begin{figure}[h]
    \centering
    \includegraphics[scale=0.8]{planner_architecture.pdf}
    \caption{Конвейер обработки запроса в PostgreSQL}
    \label{fig:postgresql_planner}
\end{figure}

При выборе плана оптимизатор опирается на статистическую информацию о данных и на параметры среды выполнения. К первой относятся сведения о размерах таблиц, распределениях значений и индексах, ко второй~--- настройки, влияющие на относительную стоимость операций ввода-вывода и вычислений, объём доступной памяти и возможности параллелизма. На основе этих данных строится стоимостная оценка альтернативных вариантов выполнения.

Процесс подготовки и выбора плана обычно включает три основных этапа: логическое переписывание запроса, генерацию альтернативных физических планов и их стоимостное сравнение. Для каждого варианта оценивается стоимость, отражающая ожидаемые затраты ресурсов или времени выполнения. Далее выбирается план с наименьшей оценочной стоимостью, который и передаётся на исполнение.

Таким образом, оптимизатор выступает как автономный модуль СУБД, скрывающий от пользователя детали физического выполнения запроса. В обычном режиме пользователь задаёт только декларативный SQL-запрос, а выбор конкретной стратегии исполнения полностью определяется оптимизатором.

\section{Проблемы и ограничения традиционных оптимизаторов}
Несмотря на многолетнее развитие, классические стоимостные оптимизаторы запросов сохраняют ряд фундаментальных ограничений \cite{Lan_2021, 10.1561/1900000002}. К основным из них относятся следующие.

\begin{itemize}
    \item \textbf{Неточность оценок кардинальностей.} \cite{10.14778/2850583.2850594, 10.1007/s00778-017-0480-7, 10.5555/645923.673638} Оценка числа строк в промежуточных результатах критически влияет на расчёт стоимости плана. В традиционных СУБД такие оценки строятся на основе упрощённых статистических предположений, прежде всего независимости предикатов и агрегированных характеристик распределений. В реальных данных эти предположения часто нарушаются, что приводит к крупным ошибкам оценки и, как следствие, к выбору неудачного плана выполнения.

    \item \textbf{Ограниченность стоимостной модели.} \cite{10.5555/645925.671662, 10.1145/1376616.1376713, 10.1109/ICDE.2013.6544899} Стоимостная модель в классических СУБД обычно выражает затраты как комбинацию вычислительной составляющей и составляющей ввода-вывода с фиксированными коэффициентами. Такая схема удобна и интерпретируема, однако она лишь приближённо отражает реальное время выполнения. На практике влияние буферного кэша, аппаратной конфигурации, параллелизма и текущего состояния системы может существенно искажать соответствие между оценочной стоимостью и фактическим временем.

    \item \textbf{Комбинаторный рост пространства планов.} \cite{10.1007/s007780050040, 10.1145/93605.98740} По мере усложнения запроса число допустимых планов быстро возрастает, поэтому даже стоимостный перебор с динамическим программированием требует отсечений и эвристик. Это позволяет удерживать время оптимизации в разумных пределах, но одновременно приводит к тому, что часть потенциально хороших альтернатив вообще не рассматривается.

    \item \textbf{Слабая адаптация к фактическому выполнению.} \cite{10.1145/276305.276315} Традиционный оптимизатор принимает решение до начала исполнения запроса и, как правило, не использует полноценную обратную связь о реальных ошибках своих оценок. Даже если во время выполнения становится очевидно, что выбранный план неудачен, СУБД обычно не способна существенно скорректировать его в текущем запуске. В результате улучшение качества планирования часто сводится к ручной перенастройке статистики, параметров или индексов.

    \item \textbf{Высокая сложность разработки и сопровождения.} \cite{10.1145/275487.275492} Оптимизатор является одним из наиболее сложных модулей СУБД. Его развитие требует одновременного согласования правил преобразования, стоимостной модели, алгоритмов перебора и механизмов поддержки новых операторов. Это делает эволюцию традиционных оптимизаторов медленной и затрудняет быстрое внедрение новых идей.
\end{itemize}

Таким образом, ограничения традиционных оптимизаторов проявляются сразу в нескольких взаимосвязанных компонентах: в оценке кардинальностей, в стоимостной модели, в алгоритмах перебора и в способности адаптироваться к изменяющейся среде выполнения. Именно эти ограничения и мотивируют развитие новых подходов, основанных на использовании статистических, обучаемых и генеративных методов\cite{2020arXiv200909884H, sel_est, 10.1145/3318464.3389741, 10.14778/3352063.3352129, 10.1145/3035918.3064029, 10.1145/3299869.3300085, marcus2019neo, ortiz2019localce, woltmann2023fastgres, anneser2023autosteer, coool2023, zhu2023lero, hero2024, ivanov2017aqo, 10.1145/3514221.3517885, 10.1145/3448016.3452838}.

\section{От стоимости к обучаемости}\label{sec:ch1/sec2}

Классический оптимизатор запросов можно рассматривать как систему, в которой качество итогового плана определяется тремя взаимосвязанными компонентами: точностью оценки кардинальностей, адекватностью стоимостной модели и эффективностью перебора пространства альтернатив. Поэтому развитие методов оптимизации запросов с использованием машинного обучения происходило не по одному направлению, а сразу по нескольким линиям. Часть работ сосредоточена на более точной оценке промежуточных объёмов данных, часть~--- на уточнении стоимостных оценок и прогнозе времени выполнения, а часть~--- на непосредственном выборе лучшего плана или набора подсказок. В последние годы к этим направлениям добавились большие языковые модели, которые используются для различных задач.

\subsection{Обучаемые методы оценки кардинальностей}

Исторически первой и наиболее изученной областью улучшения оптимизаторов стала оценка кардинальностей. Классические методы этого класса строятся на статистическом описании данных. Ранние оптимизаторы использовали простые предположения о равномерности распределений и независимости предикатов, позднее были предложены одномерные и многомерные гистограммы, а также адаптивные структуры, уточняющие статистику по мере накопления информации о рабочей нагрузке. В эту линию укладываются как классические методы на основе равновысотных и многомерных гистограмм, так и адаптивные подходы, подобные STHoles \cite{10.1145/376284.375686}, где статистическое описание уточняется на основе уже выполненных запросов.

С появлением методов машинного обучения акцент сместился от ручного построения синопсисов к обучению моделей, непосредственно предсказывающих кардинальность запроса или его подвыражений. Одной из ранних и наиболее известных работ этого направления стала MSCN \cite{kipf2018learnedcardinalitiesestimatingcorrelated}, в которой оценка кардинальности формулируется как задача регрессии по признакам таблиц, предикатов и соединений. В последующих работах стали использоваться более выразительные вероятностные и генеративные модели, например NeuroCard \cite{10.14778/3421424.3421432} и FACE \cite{10.14778/3485450.3485458}, стремящиеся приблизить совместное распределение значений в базе данных без предположений о независимости атрибутов и таблиц. Эти методы существенно повысили точность оценок на стандартных бенчмарках, особенно для сложных запросов с коррелированными условиями.

Вместе с тем литература показывает, что сама по себе высокая точность оценки кардинальностей ещё не решает всей задачи оптимизации. Во-первых, обучаемые модели требуют затрат на подготовку обучающего набора, обновление при изменении данных и контроль устойчивости к сдвигу нагрузки. Во-вторых, даже точные кардинальности не гарантируют выбора лучшего плана при несовершенной стоимостной модели или ограниченном поиске. По этой причине, хотя оценка кардинальностей остаётся центральной темой исследований, в настоящей работе она рассматривается прежде всего как фон, необходимый для понимания более прикладных направлений, связанных со стоимостной моделью, переносом подсказок и управляемым поиском.

\subsection{Обучаемые стоимостные модели}

Если оценка кардинальностей отвечает на вопрос о том, сколько строк будет обработано на каждом этапе плана, то стоимостная модель отвечает на вопрос о том, во что эти операции обойдутся с точки зрения времени и ресурсов. На практике именно здесь возникает второй крупный источник ошибок. Даже при корректной структуре плана оптимизатор может выбрать неудачную альтернативу из-за устаревших или слишком грубых параметров стоимости, не отражающих особенности конкретного оборудования, состояния буферного кэша, типа носителей и текущей нагрузки.

Поэтому отдельное направление исследований связано не с заменой оптимизатора целиком, а с построением обучаемых или гибридных стоимостных моделей. В работах \cite{10.1145/3318464.3380584} и \cite{10.1145/3511808.3557305} показывается, что время выполнения плана или отдельного оператора можно предсказывать с помощью обучаемых моделей точнее, чем это делает встроенная вручную заданная стоимостная схема. В более поздних исследованиях, например \cite{10.14778/3750601.3750645}, обсуждаются уже не только точность таких моделей, но и вопросы интерпретируемости, переносимости и внедрения в реальные системы.

При этом для практических СУБД особенно важен не столько полный отказ от классической стоимостной модели, сколько её адаптация под наблюдаемую среду выполнения. В этом случае обучаемый компонент используется как механизм уточнения или калибровки, а не как самостоятельный генератор планов. Такой подход лучше согласуется с архитектурой существующих СУБД, поскольку позволяет сохранить штатный оптимизатор и накопленные в нём эвристики, но уменьшить рассогласование между оценочной стоимостью и фактическим временем выполнения. Именно к этой линии относится и дальнейший материал диссертации: в главе 2 рассматривается не построение нового оптимизатора, а адаптивная корректировка стоимостной модели по телеметрии реальных запусков.

\subsection{Методы выбора плана и подбора подсказок}

Третье направление связано с непосредственным выбором плана или с управлением выбором через подсказки оптимизатору. В наиболее радикальной постановке обучаемая система пытается заменить классический механизм поиска полностью, то есть сама строит или ранжирует планы без опоры на традиционный оптимизатор. Однако на практике более жизнеспособными оказались методы, которые сохраняют штатный оптимизатор, но направляют его через ограниченное множество управляющих воздействий, например через подсказки или через выбор одного из нескольких кандидатных вариантов.

Наиболее известной работой этой линии является \cite{10.1145/3448016.3452838}. В ней выбор конфигурации подсказок формулируется как задача контекстного многорукого бандита, а обучаемая модель используется для того, чтобы направлять традиционный оптимизатор к более удачным планам. В дальнейшем эта линия получила развитие в работах \cite{anneser2023autosteer}, \cite{woltmann2023fastgres} и \cite{zhu2023lero}. Общая идея этих исследований состоит в том, что практическая ценность достигается не за счёт полной замены оптимизатора, а за счёт внешнего компонента, который выбирает подсказки, ранжирует планы или исправляет систематические ошибки выбора.

Для диссертации этот класс работ особенно важен по двум причинам. Во-первых, он подтверждает, что подсказки могут рассматриваться как реалистичный и технологически удобный интерфейс управления оптимизатором. Во-вторых, он показывает, что задача выбора подсказок естественно формулируется как задача поиска в дискретном пространстве альтернатив при наличии обратной связи от СУБД. Именно в этом контексте далее рассматриваются два собственных направления работы: рекомендация подсказок оптимизатору на основе сопоставления планов выполнения и агентный поиск подсказок при ограниченном бюджете пробных запусков.

\subsection{Генеративные подходы и большие языковые модели в оптимизации запросов}

Новейший этап развития связан с использованием генеративных моделей и больших языковых моделей. В отличие от более ранних обучаемых методов, ориентированных в основном на числовые признаки, такие модели позволяют работать с текстовыми и полуструктурированными представлениями запросов, планов и правил преобразования. Это расширяет спектр возможных применений: от переписывания SQL-запросов до поиска аналогий между планами и построения агентных контуров с обратной связью.

Одна группа работ использует большие языковые модели для переписывания запросов. К ней относятся, например, LLM-R2 \cite{li2024llmr2} и QUITE \cite{song2026quitequeryrewriterules}. В этих работах языковая модель используется для предложения новых вариантов переписывания, которые затем проверяются на эквивалентность и эффективность. Такой подход показывает, что генеративные модели могут выступать не только как средство генерации текста, но и как инструмент поиска новых преобразований запроса.

Другая группа работ применяет большие языковые модели к собственно задаче управления оптимизатором. В работе \cite{akioyamen2024unreasonableeffectivenessllmsquery} показано, что векторные представления текста SQL-запросов могут нести полезную информацию для выбора подсказок. 

Несмотря на высокий интерес к таким методам, литература подчёркивает и их ограничения. К ним относятся высокая стоимость генерации, нестабильность выводов, необходимость внешней валидации результатов, а также риск деградации при прямом использовании рекомендаций без защитных механизмов. Поэтому наиболее перспективными на данном этапе выглядят не сценарии полной замены оптимизатора генеративной моделью, а гибридные схемы, в которых большая языковая модель используется как внешний интеллектуальный модуль. В настоящей диссертации этот принцип реализуется в двух формах. В главе 3 большая языковая модель используется для построения векторных представлений планов и для рекомендации подсказок оптимизатору на основе сопоставления планов выполнения. В главе 4 большая языковая модель используется как агент, последовательно предлагающий подсказки и уточняющий решения по обратной связи от СУБД.

\section{Тестирование методов оптимизации запросов}\label{sec:ch1/sec3}

Корректная оценка методов оптимизации запросов требует использования воспроизводимых тестовых нагрузок, позволяющих сравнивать как итоговое время выполнения, так и качество выбора планов. В отличие от задач, связанных преимущественно с производительностью исполнения, в исследованиях оптимизаторов особенно важно, чтобы тестовый набор содержал запросы, чувствительные к ошибкам оптимизатора. К таким ошибкам относятся неточности в оценке кардинальностей, рассогласование между стоимостной моделью и фактическим временем выполнения, а также ошибки выбора порядка соединений и физических операторов.

С этой точки зрения тестовые нагрузки для методов оптимизации запросов должны оцениваться не только по популярности, но и по их способности выявлять ошибки оптимизатора. Наиболее важными характеристиками являются тип нагрузки, сложность схемы данных, разнообразие запросов, число соединений и пригодность бенчмарка для воспроизводимого тестирования. В литературе для этих целей широко используются аналитические наборы данных, поскольку именно в аналитических запросах ошибка выбора плана может приводить к значительному росту времени выполнения.

Для целей настоящей диссертации используются три бенчмарка, каждый из которых решает собственную исследовательскую задачу. Их сравнительная характеристика приведена в таблице~\ref{tab:benchmarks_selected}.

\begin{table}[h]
\centering
\small
\caption{Выбранные бенчмарки и их роль в исследовании}
\label{tab:benchmarks_selected}
\begin{tblr}{c X[c] X[c]}
\hline
\textbf{Бенчмарк} & \textbf{Основные характеристики} & \textbf{Роль в диссертации} \\
\hline
TPC-H & 
22 аналитических запроса, воспроизводимая синтетическая схема, умеренная сложность соединений &
Анализ согласованности стоимостной модели с фактическим временем выполнения, исследование влияния калибровки стоимости \\
\hline
JOB & 
Запросы по схеме IMDb, сложные и нерегулярные графы соединений, высокая чувствительность к ошибкам оптимизатора &
Стресс-тестирование качества методов рекомендации и поиска подсказок \\
\hline
CEB & 
Расширенный корпус запросов на основе JOB, существенно больший объём рабочей нагрузки &
Расширенное тестирование, оценка методов, использующих накопленный опыт, и эксперименты с обучаемыми и генеративными компонентами \\
\hline
\end{tblr}
\end{table}

Бенчмарк TPC-H используется в работе как классическая аналитическая нагрузка с хорошо воспроизводимой структурой. По сравнению с более сложными современными наборами данных он содержит относительно умеренное число шаблонов запросов и менее разнообразные структуры соединений, однако именно это делает его удобным для анализа связи между оценочной стоимостью и фактическим временем выполнения. В контексте настоящей диссертации TPC-H применяется прежде всего для исследования стоимостной модели и оценки того, насколько корректировка её параметров улучшает качество планирования.

Бенчмарк JOB, напротив, ориентирован именно на сложные случаи планирования. Он построен на реальных данных IMDb и включает запросы, в которых стандартные предположения оптимизатора о независимости предикатов и равномерности распределений часто нарушаются. Это делает JOB особенно полезным для анализа ошибок выбора порядка соединений, оценки устойчивости методов к сложным фильтрам и проверки того, действительно ли предлагаемый метод способен улучшать выбор плана, а не только подстраиваться под простые синтетические сценарии.

Наконец, CEB можно рассматривать как развитие того же направления тестирования, но уже в значительно более крупном масштабе. Если JOB удобен как компактный и хорошо интерпретируемый стресс-тест, то CEB позволяет исследовать методы на существенно большем корпусе запросов. Это особенно важно для подходов, использующих накопленный опыт выполнения, обучение на историях запросов, сопоставление планов и генеративные модели, поскольку такие методы требуют более представительной рабочей нагрузки.

Все используемые в диссертации бенчмарки относятся к классу OLAP-нагрузок. В таком классе задач критично именно качество выбора плана: большое число соединений формирует широкое пространство альтернатив, и ошибка планирования может приводить к кратному увеличению времени выполнения. Тестирование на OLTP-нагрузках в работе намеренно не проводилось, поскольку в типичных OLTP-сценариях запросы короче, планы проще, а доминирующее влияние на производительность чаще оказывает сам модуль исполнения, а не качество работы оптимизатора.

Таким образом, выбранные бенчмарки в совокупности покрывают основные потребности исследования. TPC-H используется для анализа согласованности между стоимостью и временем выполнения и для тестирования методов адаптации стоимостной модели. JOB служит основной площадкой для проверки качества выбора плана и эффективности методов рекомендации и итеративного поиска подсказок. CEB обеспечивает экспериментальную базу для тестирования методов, использующих ранее накопленный опыт, сопоставление планов и генеративные подходы. Тем самым выбранный набор тестовых нагрузок соответствует общей цели диссертации, связанной с исследованием методов оптимизации запросов на основе адаптации стоимостной модели, сопоставления планов и применения больших языковых моделей.

\section{Выводы по главе}

Проведённый обзор показывает, что классическая архитектура оптимизации SQL-запросов, сформировавшаяся на основе идей System R, по-прежнему остаётся основой современных СУБД. Её сильной стороной является универсальность: оптимизатор опирается на правила эквивалентных преобразований, стоимостную модель и алгоритмы перебора альтернативных планов, что позволяет строить решения для широкого класса запросов. Вместе с тем обзор литературы показывает и устойчивые ограничения такого подхода. К ним относятся ошибки оценки кардинальностей, приближённый характер стоимостной модели, необходимость жёстких эвристик при переборе пространства планов, а также слабая адаптация к изменяющимся данным, нагрузке и вычислительной среде.

Анализ современных исследований показывает, что развитие методов оптимизации идёт не только по пути полной замены классического оптимизатора обучаемыми моделями, но и по пути более практичных гибридных решений, совместимых с существующей архитектурой СУБД. В этих решениях машинное обучение и генеративные подходы используются для улучшения отдельных компонентов процесса оптимизации: уточнения стоимостных оценок, выбора перспективных альтернатив и переноса ранее накопленного опыта на новые запросы. Такой подход особенно важен в прикладных сценариях, где требуется не абстрактно построить новый оптимизатор, а повысить качество выбора плана в уже существующей системе.

С учётом этого в диссертации рассматриваются три взаимосвязанных направления. Во-первых, настройка стоимостной модели на основе наблюдаемой телеметрии выполняемых запросов без изменения общей архитектуры оптимизатора. Во-вторых, метод рекомендации подсказок оптимизатору на основе сопоставления планов выполнения, что позволяет использовать ранее накопленный опыт без полного перебора допустимых вариантов. В-третьих, применение языковой модели в качестве агента, который последовательно ищет эффективный набор подсказок в условиях ограниченного бюджета пробных запусков. Таким образом, работа сосредоточена не на построении универсального обучаемого оптимизатора «с нуля», а на разработке совместимых с реальными СУБД методов, улучшающих качество планирования за счёт адаптации, переноса опыта и управляемого поиска.

Отдельно проведённый обзор бенчмарков показывает, что для исследования таких методов наибольший интерес представляют аналитические нагрузки класса OLAP. Именно в них ошибка выбора плана особенно критична из-за большого числа операций соединения, сложных структур запросов и высокой чувствительности времени выполнения к качеству стоимостных оценок. По этой причине дальнейшее исследование опирается на бенчмарки TPC-H, JOB и CEB, которые в совокупности позволяют оценивать как согласованность стоимостной модели с реальным временем выполнения, так и качество выбора планов и устойчивость методов на сложных запросах.

Тем самым обзор области исследований формирует основу для последующих глав диссертации. В главе 2 рассматривается метод адаптивной настройки параметров стоимостной модели. В главе 3 исследуется метод рекомендации подсказок оптимизатору на основе сопоставления планов выполнения. В главе 4 рассматривается агентный метод итеративного поиска подсказок оптимизатору на основе большой языковой модели. Совокупно эти главы охватывают три практических способа улучшения планирования запросов без отказа от существующего оптимизатора СУБД.

\FloatBarrier
